\chapter{Brownian motion from random walks}
\label{ch:bm-walks}

\begin{question}
Where does the Wiener process come from, and why are its increments normal with variance equal to the elapsed time?
\end{question}

\begin{whatfor}
\begin{itemize}
  \item Builds Brownian motion as the limit of a fair random walk with tiny steps, and finds the one scaling that works: step size $\sigma\sqrt{\Delta t}$.
  \item Gives the density of $W(t)$ and shows that it solves the heat (diffusion) equation.
  \item Defines Brownian motion with variance parameter $\sigma^2$.
\end{itemize}
\skipline{fs:bm-walks}
\end{whatfor}

\begin{prereq}
From Part~0: normal distribution, density, mean and variance of sums, independence. From \cref{ch:first-mg}: Wiener process, random walk, martingales $W_t$ and $W_t^2-t$.
\end{prereq}

%=====================================================================
\section{A random walk with small steps}

Picture a particle on a line. Every $\Delta t$ units of time it jumps a distance $\Delta x$ to the right or to the left, with probability $\frac12$ each, independently of the past.

\begin{setup}[The scaled random walk]
Let $X_1,X_2,\ldots$ be independent with $P(X_i=+1)=P(X_i=-1)=\frac12$ ($+1$ = step right). The position at time $t$ is
\[
  B(t)=\Delta x\,\big(X_1+X_2+\dots+X_{[t/\Delta t]}\big),\qquad B(0)=0,
\]
where $[t/\Delta t]$ is the number of whole steps made by time $t$.
\end{setup}

Each step has $\E_P(X_i)=0$ and $\Var(X_i)=\E_PX_i^2=1$. By independence,
\[
  \E_P\,B(t)=0,\qquad \Var\big(B(t)\big)=(\Delta x)^2\Big[\frac{t}{\Delta t}\Big].
\]

\paragraph{Choosing the step size.} Now let $\Delta t\to0$. The variance $(\Delta x)^2\,t/\Delta t$ decides what happens:
\begin{itemize}
  \item $\Delta x=\Delta t$: the variance is $t\,\Delta t\to0$; the particle stops moving.
  \item $\Delta x=\sqrt{\Delta t}$\,: the variance is exactly $t$ in the limit; a sensible random motion.
  \item $\Delta x\gg\sqrt{\Delta t}$: the variance blows up.
\end{itemize}
So we take $\Delta x=\sigma\sqrt{\Delta t}$ for a constant $\sigma>0$, giving $\Var\,B(t)\to\sigma^2t$.

\begin{pic}
\begin{tikzpicture}
\begin{axis}[width=11cm,height=5cm,xmin=0,xmax=1,xlabel={$t$},ylabel={$B(t)$},legend style={font=\scriptsize,at={(0.5,1.02)},anchor=south},legend columns=3]
  \addplot[trapc,thick,const plot] table[x=t,y=x] {figs/scaled10.dat}; \addlegendentry{$\Delta t=\frac1{10}$}
  \addplot[thmc,thick,const plot] table[x=t,y=x] {figs/scaled100.dat}; \addlegendentry{$\Delta t=\frac1{100}$}
  \addplot[defc,thick] table[x=t,y=x] {figs/scaled1000.dat}; \addlegendentry{$\Delta t=\frac1{1000}$}
\end{axis}
\end{tikzpicture}
\picnote{Random walks with step $\Delta x=\sqrt{\Delta t}$. As the steps shrink, the staircase becomes a continuous, very rough curve: a Brownian path. In each case $\Var B(1)\approx1$.}
\end{pic}

\paragraph{The limit is normal.} For fixed $t$, $B(t)$ is $\Delta x$ times a sum of $N=[t/\Delta t]$ independent, identical steps. The central limit theorem, which we use without proof, says that a sum of $N$ independent, identically distributed steps with mean $0$ and variance $1$, divided by $\sqrt N$, has a distribution that approaches $N(0,1)$ as $N\to\infty$. Since $B(t)=\sigma\sqrt{\Delta t}\sqrt N\cdot\frac{X_1+\dots+X_N}{\sqrt N}$ and $\sqrt{\Delta t\,N}\to\sqrt t$,
\[
  B(t)\;\longrightarrow\;\sigma\sqrt t\cdot N(0,1)=N(0,\sigma^2t).
\]

The other properties pass to the limit too:
\begin{itemize}
  \item $B(0)=0$;
  \item steps in disjoint time intervals are different coins, so the increments over disjoint intervals are independent;
  \item the increment over $(s,t]$ involves about $(t-s)/\Delta t$ steps, so its law depends only on $t-s$ and tends to $N(0,\sigma^2(t-s))$.
\end{itemize}
These are the rules of the next definition. With $\sigma=1$ they are the rules of the Wiener process (\cref{ch:first-mg}).

\begin{definition}[Brownian motion with variance parameter $\sigma^2$]
\label{def:bm-sigma}
A process $(B(t))_{t\ge0}$ is a \dfn{Brownian motion with variance parameter} $\sigma^2$ if (1) $B(0)=0$; (2) its paths are continuous a.s.; (3) it has stationary and independent increments; (4) $B(t)\sim N(0,\sigma^2t)$ for every $t>0$. Then $W(t)=B(t)/\sigma$ is a standard Brownian motion, and $B(t)=\sigma W(t)$.
\end{definition}

\begin{remark}
Two more facts about Brownian paths, which we state without proof: they are continuous but nowhere differentiable (they wiggle at every scale, as the picture suggests), and a continuous process $W$ with $W(0)=0$ is a standard Brownian motion if and only if both $W_t$ and $W_t^2-t$ are martingales (\dfn{L\'evy's characterisation}).
\end{remark}

%=====================================================================
\section{The density of $W(t)$ and the heat equation}

Since $W(t)\sim N(0,t)$, its density is
\[
  f(x,t)=f_{W(t)}(x)=\frac{1}{\sqrt{2\pi t}}\,e^{-x^2/(2t)},\qquad x\in\R,\ t>0 .
\]

\begin{pic}
\begin{tikzpicture}
\begin{axis}[width=10cm,height=4.6cm,xmin=-4,xmax=4,ymin=0,ymax=0.85,xlabel={$x$},ylabel={$f(x,t)$},samples=120,legend style={font=\scriptsize}]
  \addplot[trapc,very thick,domain=-4:4] {exp(-x^2/(2*0.25))/sqrt(2*pi*0.25)}; \addlegendentry{$t=0.25$}
  \addplot[thmc,very thick,domain=-4:4] {exp(-x^2/2)/sqrt(2*pi)}; \addlegendentry{$t=1$}
  \addplot[defc,very thick,domain=-4:4] {exp(-x^2/(2*3))/sqrt(2*pi*3)}; \addlegendentry{$t=3$}
\end{axis}
\end{tikzpicture}
\picnote{The density of $W(t)$ spreads out like heat from a point: at time $t$ its standard deviation is $\sqrt t$, and the total area stays $1$.}
\end{pic}

\begin{result}[The heat equation]
\label{res:heat}
The density $f(x,t)$ of $W(t)$ satisfies the \dfn{heat equation} (diffusion equation)
\[
  \frac{\partial f}{\partial t}=\frac12\,\frac{\partial^2f}{\partial x^2}.
\]
\end{result}
\begin{proof}
Differentiate both sides and compare.
\begin{steps}
  \item Write $f=\frac{1}{\sqrt{2\pi}}t^{-1/2}e^{-x^2/(2t)}$. In $t$ (product rule):
  \[
    \frac{\partial f}{\partial t}=\frac{1}{\sqrt{2\pi}}\Big(-\frac12t^{-3/2}+t^{-1/2}\frac{x^2}{2t^2}\Big)e^{-x^2/(2t)}
    =\frac{1}{2\sqrt{2\pi}}\Big(\frac{x^2}{t^{5/2}}-\frac{1}{t^{3/2}}\Big)e^{-x^2/(2t)} .
  \]
  \item In $x$: $\dfrac{\partial f}{\partial x}=-\dfrac{x}{t}f$, and then
  \[
    \frac{\partial^2f}{\partial x^2}=-\frac1tf+\frac{x^2}{t^2}f=\frac{1}{\sqrt{2\pi}}\Big(\frac{x^2}{t^{5/2}}-\frac{1}{t^{3/2}}\Big)e^{-x^2/(2t)} .
  \]
  \item The first is half the second. \qedhere
\end{steps}
\end{proof}

\begin{intuition}
One path is rough and unpredictable. But the cloud of all paths spreads out smoothly and predictably, exactly like heat along a rod.
\end{intuition}

\begin{example}[A probability for Brownian motion with variance parameter]
$B(t)$ has variance parameter $\sigma^2=4$. Then $B(3)\sim N(0,12)$, and
\[
  P\big(B(3)>2\big)=1-\Phi\Big(\frac{2}{\sqrt{12}}\Big)=1-\Phi(0.577)\approx0.282 .
\]
\end{example}

%=====================================================================
\formulasheet{fs:bm-walks}

\begin{mastercard}[label=card:scaledwalk]{A random walk with small steps becomes Brownian motion}
Every $\Delta t$ the walk moves up $\Delta x$ with probability $p$ or down $\Delta x$ with probability $1-p$, independently. After the $t/\Delta t$ steps made by time $t$,
\[
  \boxed{\begin{gathered}\E_P\,X(t)=\frac{t}{\Delta t}\,\Delta x\,(2p-1),\qquad \Var X(t)=\frac{t}{\Delta t}\,(\Delta x)^2\big(1-(2p-1)^2\big);\\
  \Delta x=\sigma\sqrt{\Delta t},\ \ p=\tfrac12\big(1+\tfrac{\mu}{\sigma}\sqrt{\Delta t}\big),\ \ \Delta t\to0:\qquad X(t)\to\mu t+\sigma W(t)\sim N(\mu t,\sigma^2t).\end{gathered}}
\]
\begin{plugrows}
  \plugrow{Fair walk, standard Brownian motion}{\pstep{Box: $\E_PX(t)=\frac t{\Delta t}\Delta x(2p-1)$, $\Var X(t)=\frac t{\Delta t}(\Delta x)^2\big(1-(2p-1)^2\big)$.}\pstep{From ``fair'': $p=\frac12$, so $2p-1=0$; take $\Delta x=\sqrt{\Delta t}$.}\pstep{Substitute: mean $0$, variance $\frac t{\Delta t}\cdot\Delta t\cdot1=t$.}}{$\E_PX(t)=0$, $\Var X(t)=t$: $X\to W$}{$\Delta x$ must be of order $\sqrt{\Delta t}$}{$\Delta t=0.01$, $\Delta x=0.1$, $t=1$: $100$ steps; $\Var=100\cdot0.01\cdot1=1$}
  \plugrow{Variance parameter $\sigma^2$}{\pstep{Box with $p=\frac12$, $\Delta x=\sigma\sqrt{\Delta t}$.}\pstep{Substitute: mean $0$, variance $\frac t{\Delta t}\sigma^2\Delta t=\sigma^2t$.}}{$X\to\sigma W$, $\Var=\sigma^2t$}{$\sigma>0$}{$\sigma=2$, $\Delta t=0.01$: $\Delta x=2(0.1)=0.2$; $\Var X(1)=100\cdot0.04=4$}
  \plugrow{With drift $\mu$}{\pstep{Box with $\Delta x=\sqrt{\Delta t}$ ($\sigma=1$), $p=\frac12(1+\mu\sqrt{\Delta t})$, so $2p-1=\mu\sqrt{\Delta t}$.}\pstep{Mean: $\frac t{\Delta t}\sqrt{\Delta t}\,\mu\sqrt{\Delta t}=\mu t$.}\pstep{Variance: $\frac t{\Delta t}\Delta t(1-\mu^2\Delta t)\to t$.}}{$X(t)\to\mu t+W(t)$}{$|\mu|\sqrt{\Delta t}\le1$}{$\mu=1$, $\Delta t=0.01$: $p=\frac12(1.1)=0.55$, $\E_PX(1)=100(0.1)(2\cdot0.55-1)=100(0.1)(0.1)=1$}
  \plugrow{Wrong scalings}{\pstep{Box with $p=\frac12$: $\Var X(t)=\frac t{\Delta t}(\Delta x)^2$.}\pstep{$\Delta x=\Delta t$: $\Var=t\,\Delta t\to0$.}\pstep{$\Delta x\gg\sqrt{\Delta t}$: $\frac{(\Delta x)^2}{\Delta t}\to\infty$.}}{the motion collapses, or blows up}{only $\Delta x\propto\sqrt{\Delta t}$ works}{$\Delta x=\Delta t=0.01$: $\Var X(1)=100\cdot0.0001=0.01$}
  \plugrow{Multiplicative walk}{\pstep{The price is multiplied by $e^{\pm\sigma\sqrt h}$ every $h$, so $\ln(\text{price})$ moves by $\pm\sigma\sqrt h$.}\pstep{Box for $\ln(\text{price})$ with $\Delta t=h$, $\Delta x=\sigma\sqrt h$, $p=\frac12(1+\frac\mu\sigma\sqrt h)$: mean $\mu t$, variance $\to\sigma^2t$.}\pstep{Exponentiate.}}{$\ln(\text{price})\to\ln(\text{start})+\mu t+\sigma W_t$, i.e.\ price $=\text{start}\cdot e^{\mu t+\sigma W_t}$ (Card 11.1)}{log-steps $\pm\sigma\sqrt h$}{$\mu=0.1$, $\sigma=0.2$: $\ln(\text{price})$ has mean $0.1t$, variance $0.2^2t=0.04t$}
  \plugrow{Density of the limit: the heat equation}{\pstep{Box, limit: $X(t)\sim N(x_0+\mu t,\sigma^2t)$, density $f(t,y)=\frac1{\sqrt{2\pi\sigma^2t}}e^{-(y-x_0-\mu t)^2/2\sigma^2t}$.}\pstep{Differentiate once in $t$ and twice in $y$.}\pstep{Compare the three derivatives.}}{$f_t=-\mu f_y+\frac{\sigma^2}{2}f_{yy}$; for $W$: $f_t=\frac12f_{yy}$}{$t>0$}{$W$ at $y=0$, $t=1$: $f=\frac1{\sqrt{2\pi}}$, $f_t=-\frac12\cdot\frac{1}{\sqrt{2\pi}}$, $f_{yy}=-\frac1{\sqrt{2\pi}}$, and $-\frac{1}{2\sqrt{2\pi}}=\frac12\big(-\frac1{\sqrt{2\pi}}\big)$}
\end{plugrows}
\end{mastercard}
